\documentclass[11pt]{article}
\newcommand{\ListaTitulo}{Gabarito 04: Correlação, assimetria e curtose}
% Preâmbulo compartilhado: MATF14, Listas de Exercícios 2026
% Usado por ListaNN.tex e GabaritoNN.tex via % Preâmbulo compartilhado: MATF14, Listas de Exercícios 2026
% Usado por ListaNN.tex e GabaritoNN.tex via % Preâmbulo compartilhado: MATF14, Listas de Exercícios 2026
% Usado por ListaNN.tex e GabaritoNN.tex via \input{preamble}

\usepackage[utf8]{inputenc}
\usepackage[T1]{fontenc}
\usepackage[brazil]{babel}
\usepackage{fullpage}
\usepackage{amsmath,amssymb}
\usepackage{enumitem}
\usepackage{xcolor}
\usepackage{fancyhdr}
\usepackage{titlesec}
\usepackage{listings}
\usepackage{hyperref}
\usepackage{booktabs}
\usepackage{graphicx}
\usepackage{tikz}

\definecolor{matf14azul}{HTML}{0D3B54}
\definecolor{matf14dourado}{HTML}{C9971E}

\titleformat{\section}{\normalfont\Large\bfseries\color{matf14azul}}{\thesection}{1em}{}
\titleformat{\subsection}{\normalfont\large\bfseries\color{matf14azul}}{\thesubsection}{1em}{}

\providecommand{\ListaTitulo}{Lista de Exercícios}

\setlength{\headheight}{14pt}
\pagestyle{fancy}
\fancyhf{}
\lhead{\small MATF14: Estatística Econômica I}
\rhead{\small \ListaTitulo}
\cfoot{\thepage}
\renewcommand{\headrulewidth}{0.4pt}

\lstset{
  basicstyle=\ttfamily\small,
  breaklines=true,
  frame=single,
  columns=fullflexible,
  backgroundcolor=\color{gray!8},
  commentstyle=\color{matf14dourado},
  keywordstyle=\color{matf14azul}\bfseries,
  language=R
}

\newcounter{questaocounter}
\newenvironment{questao}[1]{%
  \stepcounter{questaocounter}%
  \bigskip\noindent\textbf{Questão #1.}\ %
}{\par}

\newcommand{\cabecalho}[2]{%
  \begin{center}
    {\Large\bfseries\color{matf14azul} #1}\\[0.3em]
    {\large #2}
  \end{center}
  \vspace{0.5em}
  \noindent\rule{\linewidth}{0.6pt}
  \vspace{0.5em}
}

\newenvironment{solucao}{%
  \par\smallskip\noindent\textit{\color{matf14dourado!80!black}Solução.}\ %
}{\hfill$\blacksquare$\par}


\usepackage[utf8]{inputenc}
\usepackage[T1]{fontenc}
\usepackage[brazil]{babel}
\usepackage{fullpage}
\usepackage{amsmath,amssymb}
\usepackage{enumitem}
\usepackage{xcolor}
\usepackage{fancyhdr}
\usepackage{titlesec}
\usepackage{listings}
\usepackage{hyperref}
\usepackage{booktabs}
\usepackage{graphicx}
\usepackage{tikz}

\definecolor{matf14azul}{HTML}{0D3B54}
\definecolor{matf14dourado}{HTML}{C9971E}

\titleformat{\section}{\normalfont\Large\bfseries\color{matf14azul}}{\thesection}{1em}{}
\titleformat{\subsection}{\normalfont\large\bfseries\color{matf14azul}}{\thesubsection}{1em}{}

\providecommand{\ListaTitulo}{Lista de Exercícios}

\setlength{\headheight}{14pt}
\pagestyle{fancy}
\fancyhf{}
\lhead{\small MATF14: Estatística Econômica I}
\rhead{\small \ListaTitulo}
\cfoot{\thepage}
\renewcommand{\headrulewidth}{0.4pt}

\lstset{
  basicstyle=\ttfamily\small,
  breaklines=true,
  frame=single,
  columns=fullflexible,
  backgroundcolor=\color{gray!8},
  commentstyle=\color{matf14dourado},
  keywordstyle=\color{matf14azul}\bfseries,
  language=R
}

\newcounter{questaocounter}
\newenvironment{questao}[1]{%
  \stepcounter{questaocounter}%
  \bigskip\noindent\textbf{Questão #1.}\ %
}{\par}

\newcommand{\cabecalho}[2]{%
  \begin{center}
    {\Large\bfseries\color{matf14azul} #1}\\[0.3em]
    {\large #2}
  \end{center}
  \vspace{0.5em}
  \noindent\rule{\linewidth}{0.6pt}
  \vspace{0.5em}
}

\newenvironment{solucao}{%
  \par\smallskip\noindent\textit{\color{matf14dourado!80!black}Solução.}\ %
}{\hfill$\blacksquare$\par}


\usepackage[utf8]{inputenc}
\usepackage[T1]{fontenc}
\usepackage[brazil]{babel}
\usepackage{fullpage}
\usepackage{amsmath,amssymb}
\usepackage{enumitem}
\usepackage{xcolor}
\usepackage{fancyhdr}
\usepackage{titlesec}
\usepackage{listings}
\usepackage{hyperref}
\usepackage{booktabs}
\usepackage{graphicx}
\usepackage{tikz}

\definecolor{matf14azul}{HTML}{0D3B54}
\definecolor{matf14dourado}{HTML}{C9971E}

\titleformat{\section}{\normalfont\Large\bfseries\color{matf14azul}}{\thesection}{1em}{}
\titleformat{\subsection}{\normalfont\large\bfseries\color{matf14azul}}{\thesubsection}{1em}{}

\providecommand{\ListaTitulo}{Lista de Exercícios}

\setlength{\headheight}{14pt}
\pagestyle{fancy}
\fancyhf{}
\lhead{\small MATF14: Estatística Econômica I}
\rhead{\small \ListaTitulo}
\cfoot{\thepage}
\renewcommand{\headrulewidth}{0.4pt}

\lstset{
  basicstyle=\ttfamily\small,
  breaklines=true,
  frame=single,
  columns=fullflexible,
  backgroundcolor=\color{gray!8},
  commentstyle=\color{matf14dourado},
  keywordstyle=\color{matf14azul}\bfseries,
  language=R
}

\newcounter{questaocounter}
\newenvironment{questao}[1]{%
  \stepcounter{questaocounter}%
  \bigskip\noindent\textbf{Questão #1.}\ %
}{\par}

\newcommand{\cabecalho}[2]{%
  \begin{center}
    {\Large\bfseries\color{matf14azul} #1}\\[0.3em]
    {\large #2}
  \end{center}
  \vspace{0.5em}
  \noindent\rule{\linewidth}{0.6pt}
  \vspace{0.5em}
}

\newenvironment{solucao}{%
  \par\smallskip\noindent\textit{\color{matf14dourado!80!black}Solução.}\ %
}{\hfill$\blacksquare$\par}


\begin{document}

\cabecalho{Gabarito 04}{Coeficiente de correlação, assimetria e curtose: fecha a Unidade 1 (Cap. 1, Aula 08)}

\begin{questao}{1} \textbf{(Interpretação de $r$: gastos com propaganda e vendas)}
\begin{solucao}
\begin{enumerate}[label=(\alph*)]
  \item Existe uma associação linear forte e positiva entre gasto em propaganda e faturamento do
    mês seguinte: lojas que gastaram mais tenderam a faturar mais no mês seguinte, e essa
    tendência é bastante consistente (próxima de uma reta).
  \item Lojas maiores (mais movimento, mais clientes de base) podem gastar mais em propaganda \emph{e}
    faturar mais, simplesmente por serem maiores, sem que a propaganda seja a causa do
    faturamento adicional. Porte da loja é uma variável de confusão plausível.
  \item Um experimento simples: sortear aleatoriamente, entre as lojas, quais recebem um aumento
    no orçamento de propaganda e quais mantêm o orçamento atual, e comparar o faturamento
    subsequente entre os dois grupos. Ao aleatorizar quem recebe o "tratamento" (mais propaganda),
    o porte da loja fica balanceado entre os grupos, permitindo isolar o efeito da propaganda.
\end{enumerate}
\end{solucao}
\end{questao}

\begin{questao}{2} \textbf{(Correlação quase nula não é ausência de relação)}
\begin{solucao}
\begin{enumerate}[label=(\alph*)]
  \item $r$ mede exclusivamente associação \textbf{linear}. Sim, isso limita a conclusão: $r\approx 0$
    só permite concluir que não há uma tendência de reta entre as duas variáveis, não permite
    concluir ausência de qualquer relação, inclusive relações fortes porém não lineares.
  \item Por exemplo, se vendas de imóveis caem quando os juros sobem \emph{ou} quando os juros caem
    demais (por sinalizar recessão/instabilidade), e sobem quando os juros estão em um patamar
    intermediário "saudável", uma relação em forma de "U invertido", os pontos altos e baixos de
    juros associados a vendas baixas cancelariam a inclinação de uma reta ajustada aos dados,
    produzindo $r$ próximo de zero mesmo com uma relação sistemática e previsível (só que curva,
    não linear).
\end{enumerate}
\end{solucao}
\end{questao}

\begin{questao}{3} \textbf{(Assimetria: comparando duas distribuições salariais)}
\begin{solucao}
\begin{enumerate}[label=(\alph*)]
  \item X: histograma aproximadamente simétrico, em formato de sino, centrado em R\$ 5.000. Y:
    histograma com um pico concentrado em valores mais baixos e uma cauda longa se estendendo
    para a direita (poucos salários muito altos), mesmo com média e desvio padrão idênticos aos
    de X.
  \item Na empresa Y (assimetria positiva forte): como vimos na Aula 04/Cap. 1, em distribuições
    assimétricas à direita a relação usual é moda $<$ mediana $<$ média, a mediana fica abaixo da
    média, puxada para baixo pela concentração de valores menores, enquanto a média é puxada para
    cima pela cauda de salários altos.
  \item Um reajuste fixo (somar a mesma constante $k$ a todos os salários) desloca toda a
    distribuição para a direita, mas \textbf{não altera sua forma}: a assimetria, que mede o
    desequilíbrio da distribuição em torno da própria média (uma quantidade adimensional,
    normalizada pelo desvio padrão ao cubo), permanece a mesma, um reajuste fixo não torna a
    distribuição salarial mais nem menos assimétrica.
\end{enumerate}
\end{solucao}
\end{questao}

\begin{questao}{4} \textbf{(Curtose e risco financeiro)}
\begin{solucao}
\begin{enumerate}[label=(\alph*)]
  \item Está \textbf{subestimando}: curtose 7 (bem acima de 3) indica caudas mais pesadas que a Normal, eventos extremos (quedas ou altas bruscas) ocorrem com frequência maior do que um modelo
    Normal preveria. Assumir Normal para Q leva a subestimar o risco de eventos extremos.
  \item Porque a curtose afeta diretamente a frequência esperada de perdas muito grandes (o que
    interessa à gestão de risco, como no cálculo de um Value-at-Risk), enquanto a média de retorno
    sozinha nada diz sobre a probabilidade desses eventos raros e extremos.
\end{enumerate}
\end{solucao}
\end{questao}

\begin{questao}{5} \textbf{(Dedução: invariância de $r$ à mudança de unidade)}
\begin{solucao}
\begin{enumerate}[label=(\alph*)]
  \item Por definição, $\overline{X'} = \frac{1}{n}\sum_{i=1}^n x_i' = \frac{1}{n}\sum_{i=1}^n
    (ax_i+b) = a\left(\frac{1}{n}\sum_i x_i\right) + b = a\bar x + b$. Logo,
    $$
    x_i' - \overline{X'} = (ax_i+b) - (a\bar x + b) = a(x_i-\bar x).
    $$
  \item Substituindo $x_i'-\overline{X'} = a(x_i-\bar x)$ na definição de $r_{X'Y}$:
    $$
    r_{X'Y} = \frac{\sum_i a(x_i-\bar x)(y_i-\bar y)}{\sqrt{\sum_i a^2(x_i-\bar x)^2}\ \sqrt{\sum_i (y_i-\bar y)^2}}
    = \frac{a\sum_i (x_i-\bar x)(y_i-\bar y)}{a\sqrt{\sum_i (x_i-\bar x)^2}\ \sqrt{\sum_i (y_i-\bar y)^2}},
    $$
    usando que $a>0 \Rightarrow \sqrt{a^2} = a$, que sai do somatório do denominador (é constante em
    $i$). Como $a>0$ aparece multiplicando tanto o numerador quanto o denominador, ele se cancela:
    $$
    r_{X'Y} = \frac{\sum_i (x_i-\bar x)(y_i-\bar y)}{\sqrt{\sum_i (x_i-\bar x)^2}\ \sqrt{\sum_i (y_i-\bar y)^2}} = r_{XY}.
    $$
    Portanto $r$ é invariante a mudanças lineares de escala com coeficiente positivo, trocar
    dólares por reais (multiplicando por uma cotação fixa positiva) não altera o valor da
    correlação medida.
\end{enumerate}
\end{solucao}
\end{questao}

\end{document}
